\section{The continuation theorem}
\label{sec:continuation}

We now put the three checks together. In words: a handoff whose certificate
passes has a target that realizes the conditions, a strict live-overlap
lineage from source to target, and a target that keeps itself going once
everything used for the transfer is withdrawn.

\subsection{The certificate}

We call a \emph{handoff model} the following data: a source and a target system, typed maps
that carry each condition's cells from source to target, a cover on the
target, the stages and routes of the episode, a record of the execution, and a
list of resources marked as withdrawn or retained. All of these are given as
raw data. The predicates that the checks use are defined from that data, and
no derived object is supplied separately: not the target's assignment, not the
predictive classes, not the success flags. The one thing that is supplied is
evidence, namely a proof that an invariance obligation holds, used by the
autonomy check. Given an assignment of values to
the source's cells, we compute the target's assignment by mapping the
source's local pieces across and gluing them on the target's cover.

\begin{definition}[Continuation certificate]
% claim: TTH-06-GEN
\label{def:certificate}
A handoff model with a source assignment passes the \emph{continuation
certificate} if every check of the formal gate holds. Appendix~\ref{app:checks}
expands the gate component by component. In summary, the gate requires the following.
\begin{enumerate}
\item \emph{Source.} The source assignment satisfies the eight conditions on
the source.
\item \emph{Realization.} The mapped pieces satisfy the cover, scope,
totality and value-preservation checks. Each is accepted by its own condition
on the target, and adjacent pieces agree where they overlap.
\item \emph{Lineage.} For every route of the episode the computed handoffs
form a strict chain (Definition~\ref{def:handoff}). The predictive and
record readouts of source and target are related by the declared transport.
Within each of the two systems, states with the same predictive code have the
same observation laws for every finite sequence of actions. The translation
from source to target commutes with initial states, transition weights,
successors and observations on the whole set of source states the execution
reaches. The typed endpoint obligations and the coherence of the shared
runtime also hold, including the temporal condition's decoder.
\item \emph{Autonomy.} After every declared role is withdrawn, a game
invariant certificate (Lemma~\ref{lem:strategy}) covers the target's runs. It
is a stronger object than the finite certificate of
Definition~\ref{def:autonomy}: it ranges over full run histories, and it
requires normalized responses in the stated sparse format, legal moves using
only the target's own retained resources, and strict advance of order, tick,
progress and records. In the \emph{realization-preserving} version the
invariant includes the target's realization. In the \emph{operational}
version it requires only membership in a declared part of the target's state
space.
\item \emph{Bookkeeping.} The route is recorded as route-indexed, and the
receiver of the source's content passes its own content and intervention
checks.
\end{enumerate}
\end{definition}

\subsection{The theorem}

\begin{theorem}[Continuation theorem]
% claim: TTH-06-GEN
\label{thm:continuation}
Let a handoff model with a source assignment pass the continuation certificate
in its realization-preserving version. Then the following hold.
\begin{enumerate}
\item[(R)] \emph{Realization.} The computed target assignment is total,
restricts to every mapped piece on its scope, and is accepted by all eight
conditions on the target.
\item[(L)] \emph{Lineage.} For every route the computed handoffs form a
strict chain. Moreover, within the source and within the target, any two
states with the same predictive code give the same probability to every finite
sequence of observations under every finite sequence of actions.
\item[(A)] \emph{Autonomy.} After withdrawal there is a strategy for the target
such that every history it generates, from every genuine entry state and
following every outcome of positive probability, stays realized, uses only the
target's own retained resources, makes only legal and affordable moves, and
strictly advances order, tick, progress and records at every step.
\end{enumerate}
In the operational version, (R) and (L) hold as stated, and (A) holds with
``stays realized'' replaced by ``stays in the declared part of the state
space''.
\end{theorem}

We prove the theorem from three lemmas, and we state the two that carry most
of the weight.

\begin{lemma}[From one step to every future]
% claim: TTH-06-GEN
\label{lem:lumping}
Consider a system with a finite list of outcomes, and suppose any two states
with the same code satisfy three conditions. They have the same observation.
For every action and every outcome, they move to states with the same code.
For every action, they give equal total weight to each code class of next
states. Then any two states with the same code give the same probability to
every finite sequence of observations, under every finite sequence of
actions.
\end{lemma}

\begin{lemma}[From an invariant to a strategy]
% claim: TTH-06-GEN
\label{lem:strategy}
Suppose an invariant set of runs satisfies the following conditions.
\begin{itemize}
\item The withdrawal holds, and some entry run exists.
\item Every entry lies in the invariant.
\item Every run in the invariant satisfies the realization predicate and the
resource conditions.
\item For every run in the invariant and every disturbance, some legal action
has normalized outcome weights (nonnegative, summing to one, with a positive
entry), and each outcome of positive weight leads to a run in the invariant
and satisfies the transition law.
\end{itemize}
Then there is a causal strategy under which every history generated from every
entry consists of runs satisfying the realization predicate and the resource
conditions. At each of these runs, for every disturbance, the chosen action is
legal and normalized, and each of its positive outcomes satisfies the
transition law.
\end{lemma}

Lemma~\ref{lem:lumping} is a lumpability argument \citep{KemenySnell1976},
carried from a single chain to a system with actions and observations, by
induction on the length of the sequence. Its hypothesis makes the partition
into codes a probabilistic bisimulation in the sense of \citet{LarsenSkou1991},
adapted so that observations label the states. It is stronger still, since it aligns
successor codes outcome by outcome and not only the masses of the code
classes. Its conclusion is that states with the same code agree on every
action-observation test in the sense of \citet{LittmanSuttonSingh2002}.
Lemma~\ref{lem:strategy} is the standard construction of a
controlled-invariant strategy \citep{Blanchini1999, AubinBayenSaintPierre2011}.
In the terms of infinite games, the invariant is a trap for the environment,
which chooses the disturbance and the outcome, and the strategy keeps play
inside it \citep[p.~36]{Mazala2002}.
Its hypotheses exclude the vacuous cases: an empty set of entries, and empty
responses. The third ingredient is the gluing theorem of
Section~\ref{sec:realization}, together with the locality of each
condition, which carries local acceptance over to the glued assignment. With
these in hand, the proof of the theorem consists of applying each lemma to its
own check. We say this plainly, because the theorem is only as strong as its
definitions. Its value lies in two facts. Each check's predicates are defined from the
model's own data. Each conclusion is a property of the model, not merely a field
of a certificate, even though autonomy draws on a supplied proof of invariance.

\paragraph{Two versions.}
In the realization-preserving version, autonomy implies realization, because
the invariant includes it. In the operational version, the three conclusions
are logically independent, in the following sense. On the two-copy model we
make three single edits of the baseline: deleting one published cell, removing
the source's support from one overlap step, or withdrawing the target's
retained recharge. Each edit makes exactly one of the three conclusions fail while the
other two hold. We do not exhibit all eight patterns of success and failure.

\paragraph{What the theorem does not give.}
We prove no converse: a target that realizes the conditions and continues the
source need not have a passing certificate. The autonomy conclusion covers the
disturbances of the recorded execution and every outcome of positive
probability of the chosen response. We say nothing about resource shocks or
perturbations of the transition weights. The sparse format in which responses
are listed is a representation choice, and the certificate depends on it.
Autonomy in our sense is equivalent to the existence of some invariant
certificate. This is a statement about existence, not a procedure for finding
one. The finite construction we use to produce certificates is sound but may
be incomplete, so its failure to produce one would not show that autonomy
fails.
Lineage is strict lineage. It is not identity, and realization on the target is
not the soundness of any certificate for consciousness. Transport of the
temporal signature alone does not establish lineage.

\subsection{Instances}

\begin{proposition}[A handoff that passes]
% claim: TTH-06-INST
\label{prop:instance}
On the two-copy model of Section~\ref{sec:copy}, with its baseline
assignment, the continuation certificate passes in both versions. Hence one
handoff from the source copy to the target copy has realization on the target,
strict lineage and autonomy after withdrawal.
\end{proposition}

The handoff is a single change of system. Conditions (4) to (7) hold in the
simplified form of Section~\ref{sec:setting}. The autonomy check draws on the
source's support in the recorded history, while the lineage check uses the
source's funded execution, and we keep the two apart. A second instance, with
the histories of the two routes exchanged, also passes.

A single handoff is one step. The second instance is a chain.

\begin{proposition}[A chain through three systems]
% claim: TTH-07
\label{prop:chain-instance}
On the chain model, an episode passes through three distinct systems, labelled
\snapGSevenStrictChainSystemIds{} in order, by two handoffs. For each of the
\snapGSevenNativeFactorStates{} states of its finite factor, the following hold:
the first and last systems are distinct; the final target owns itself; both
handoffs realize the conditions on their targets; both satisfy their lineage
checks; and the final target's controller wins its game after withdrawal. The
middle system's whole signature agrees exactly between the two handoffs.
\end{proposition}

The composite lineage here is the conjunction of the two handoffs' lineages.
The general fact that strict handoffs compose into a strict chain is
Proposition~\ref{prop:chain}. On the chain model, conditions (4) to (7) again
hold in simplified form, and disturbances are limited to a unit class. Several
of the withdrawn roles coincide in this model. Its receiver preserves the
payload and the source's classes, but not which history or route the source
took. A general composition statement for such certificates also holds, but it is a
formal rearrangement that does not use the agreement between the two
handoffs, so we rely on the checked instance instead.
Figure~\ref{fig:chain} draws the chain.

\begin{figure}[tbp]
\centering
% F4: distinct system IDs and refinement are selectors of pinned snapshot lists.
\begin{tikzpicture}[x=1cm,y=1cm,font=\small,
 system/.style={draw,rounded corners=2pt,fill=black!3,minimum width=3cm,minimum height=.95cm,align=center},
 flow/.style={-{Stealth[length=2mm]},semithick}]
\node[system] (source) at (1.7,2.3) {Source system\\\snapGSevenStrictChainSystemIdsItemZero};
\node[system] (middle) at (7.1,2.3) {Intermediate system\\\snapGSevenStrictChainSystemIdsItemOne};
\node[system] (target) at (12.5,2.3) {Final target system\\\snapGSevenStrictChainSystemIdsItemTwo};
\draw[flow] (source)--node[above,align=center] {live\\handoff} (middle);
\draw[flow,densely dashed] (middle)--node[above,align=center] {live\\handoff} (target);
\node[draw,densely dotted,align=center,text width=3cm,minimum height=.9cm] (first) at (4.4,.25)
 {Source withdrawn;\\intermediate continues};
\node[draw,densely dotted,align=center,text width=3cm,minimum height=.9cm] (second) at (9.8,.25)
 {Intermediate withdrawn;\\target continues};
\draw[flow,densely dotted] (4.4,2.3)--(first);
\draw[flow,densely dotted] (9.8,2.3)--(second);
\node[draw,rounded corners=2pt,align=center,text width=12.7cm] at (7.1,-1.3)
 {Predictive refinement: \snapGSevenCThreeTrajectoryItemZero{} $\longrightarrow$
 \snapGSevenCThreeTrajectoryItemOne{} $\longrightarrow$ \snapGSevenCThreeTrajectoryItemTwo{}
 \quad (stabilized)\\Non-exposure, action credit, temporal closure and binding are simplified};
\end{tikzpicture}

\caption{The chain through three systems. Each handoff is live, and each earlier system is withdrawn after its handoff. The box gives the predictive refinement on the model's finite factor. Conditions (4) to (7) hold in simplified form on this model.}
\label{fig:chain}
\end{figure}
